% ch10_b 第10章 §10.2尾–§10.3 (书页 453–465 / p468–p480)
% 块界说明：
%   起点——书页 453（p468）顶部 "will be called a string-net operator." 是书页 452
%   （p467，ch10_a 块）末行显示公式 W(C_net)=W(C_1)W(C_2)… 之后半句的延续；
%   ch10_a 已把前半句译文（"……将被称为"）留在其正文最后。本块正文以 %JOIN%
%   起头续译"弦网算符。……"，不设 %--A-TAIL-- 区，亦不重复前半句。
%   终点——书页 465（p480）末句 "The PSGs obtained in Section 10.2 ... agree
%   exactly with the" 延续到书页 466（p481，ch10_c 块）顶部 "fermion PSGs that
%   we obtained above. ..."。按约定，半句译文（"……与"）留在本文件正文最末，
%   未写入任何 TAIL 区；请 ch10_c 以 %JOIN% 起头续译"我们在上面得到的费米子
%   PSG 完全一致。……"，勿重复本块已有内容。
%JOIN%
弦网算符（string-net operator）。只要 $C_i$ 中至少有一个是开弦，态 $|C_1C_2\dots\rangle$ 就是开弦网态。相应的算符 $W(C_{\mathrm{net}})$ 将被称为开弦网算符。若所有 $C_i$ 都是闭合回路，则 $|C_1C_2\dots\rangle$ 就是闭合弦网态，而 $W(C_{\mathrm{net}})$ 是闭合弦网算符。

哈密顿量 $H_J$ 没有弦网凝聚，因为它的基态 $|0\rangle$ 不包含任何弦网，因而 $\langle0|W(C_{\mathrm{net}})|0\rangle=0$。要得到具有闭合弦网凝聚的哈密顿量，我们首先需要找到一个哈密顿量，其基态包含大量任意尺寸的闭合弦网，并且不包含任何开弦（或开弦网）。

让我们先写出一个使闭合弦和闭合弦网不耗费能量、而任何开弦都耗费大量能量的哈密顿量。这样一个哈密顿量可以取如下形式
\begin{equation*}
H_U=-U\sum_{\mathrm{even}}\hat F_{\pmb{i}}
\end{equation*}
其中
\begin{equation*}
\hat F_{\pmb{i}}=\sigma_{\pmb{i}}^{x}\sigma_{\pmb{i}+\pmb{x}}^{y}\sigma_{\pmb{i}+\pmb{x}+\pmb{y}}^{x}\sigma_{\pmb{i}+\pmb{y}}^{y}\tag{10.3.2}
\end{equation*}

我们发现，无弦态 $|0\rangle$ 是 $H_U$ 的基态之一（假定 $U>0$），能量为 $-UN_{\mathrm{site}}$。所有闭合弦态，例如 $W(C_{\mathrm{close}})|0\rangle$，也都是 $H_U$ 的基态，因为 $[H_U,W(C_{\mathrm{close}})]=0$。类似地，$[H_U,W(C_{\mathrm{net}})]=0$ 意味着任何闭合弦态也都是基态。

利用开弦算符 $W(C_{\mathrm{open}})$ 与 $\hat F_{\pmb{i}}$ 之间的对易关系
\begin{align*}
\hat F_{\pmb{i}}W(C_{\mathrm{open}})&=-W(C_{\mathrm{open}})\hat F_{\pmb{i}}, && \text{若 } C_{\mathrm{open}} \text{ 终止于方格 } \pmb{i},\\
\hat F_{\pmb{i}}W(C_{\mathrm{open}})&=+W(C_{\mathrm{open}})\hat F_{\pmb{i}}, && \text{其他情形},
\end{align*}
我们发现，开弦算符在其两个端点处翻转 $\hat F_{\pmb{i}}$ 的符号。开弦态 $W(C_{\mathrm{open}})|0\rangle$ 也是 $\hat F_{\pmb{i}}$ 的本征态，因而是 $H_U$ 的本征态，能量为 $-UN_{\mathrm{site}}+4U$。我们看到，开弦的每个端点耗费能量 $2U$。我们还注意到，闭合弦网的能量不依赖于网中弦的长度。因此，$H_U$ 中的弦没有张力。可以通过在哈密顿量中加入 $H_J$ 来引入弦张力，张力为每格点（或每段）$2J$。我们注意到，任何弦网态 $|C_1C_2\dots\rangle$ 都是 $H_U+H_J$ 的本征态。因此，由 $H_U+H_J$ 描述的弦网不涨落，从而不能凝聚。为使弦发生涨落，我们需要 $g$ 项
\begin{equation*}
H_g=-g\sum_{\pmb{p}}W(C_{\pmb{p}})=-g\sum_{\mathrm{odd}}\hat F_{\pmb{i}}
\end{equation*}
其中 $\pmb{p}$ 标记奇数方格，$C_{\pmb{p}}$ 是围绕方格 $\pmb{p}$ 的闭合弦。当 $H_g$ 作用在弦网态上时，它往弦网中添加一圈弦 $C_{\pmb{p}}$。因此 $g$ 项使弦发生涨落。我们的自旋 1/2 模型的总哈密顿量为
\begin{equation*}
H=H_U+H_J+H_g
\end{equation*}

当 $U\gg g\gg J$ 时，闭合弦几乎没有张力，$H_g$ 诱发的涨落不受张力的限制。因此，基态包含闭合弦的强烈涨落。换言之，基态被任意尺寸的闭合弦所充满。若 $U\gg g$，则闭合弦断开成开弦要耗费过多能量。因此，当 $U\gg g\gg J$ 时，我们自旋 1/2 模型的基态包含闭合弦凝聚（或者更确切地说，闭合弦网凝聚）。下面我们将研究这样一种弦网凝聚态的物理性质。

\subsection*{问题 10.3.1}

(a) 证明各 $\hat F_{\pmb{i}}$ 在 $x$ 与 $y$ 两个方向均为周期性边界条件的格子上满足
\begin{equation*}
\prod_{\pmb{i}}\hat F_{\pmb{i}}=1\tag{10.3.3}
\end{equation*}
式 (10.3.2) 中的负号使上述简单结果成为可能。

(b) 在两方向均为周期性边界条件的偶 $\times$ 偶格子上，证明各 $\hat F_{\pmb{i}}$ 满足更强的条件
\begin{equation*}
\prod_{\pmb{i}=\mathrm{even}}\hat F_{\pmb{i}}=1\tag{10.3.4}
\end{equation*}

\subsection{弦网凝聚与低能有效理论}

\begin{keypoints}\item 弦网凝聚给出规范激发。\end{keypoints}

让我们先讨论 $U\gg g>0$、$J=0$ 的弦网凝聚相。当 $J=0$ 时，模型 $H_U+H_g$ 成为上一节讨论过的模型 (10.2.6)，其中在偶数格点 $V_{\pmb{i}}=U$，在奇数格点 $V_{\pmb{i}}=g$。该模型是精确可解的，因为 $[\hat F_{\pmb{i}},\hat F_{\pmb{j}}]=0$（Kitaev, 2003）。$H_U+H_g$ 的本征态可以由 $\hat F_{\pmb{i}}$ 的共同本征态得到，即 $|\{F_{\pmb{i}}\}\rangle$，其中 $F_{\pmb{i}}$ 是 $\hat F_{\pmb{i}}$ 的本征值。由于 $\hat F_{\pmb{i}}^{2}=1$，$\hat F_{\pmb{i}}$ 的本征值就是 $F_{\pmb{i}}=\pm1$。$|\{F_{\pmb{i}}\}\rangle$ 的能量为 $E=-U\sum_{\pmb{i}=\mathrm{even}}F_{\pmb{i}}-g\sum_{\pmb{i}=\mathrm{odd}}F_{\pmb{i}}$。

注意到，对于尺寸为 $L_x\times L_y$ 的偶 $\times$ 偶周期格子上的有限系统，各 $\hat F_{\pmb{i}}$ 满足约束 $\prod_{\pmb{i}=\mathrm{even}}\hat F_{\pmb{i}}=1$ 与 $\prod_{\pmb{i}=\mathrm{odd}}\hat F_{\pmb{i}}=1$。因此，各 $F_{\pmb{i}}$ 并不独立，只能标记 $2^{L_xL_y}/4$ 种不同的组态。然而，我们自旋 1/2 模型的希尔伯特空间包含 $2^{L_xL_y}$ 个态。为了复现这 $2^{L_xL_y}$ 个态，$\hat F_{\pmb{i}}$ 的共同本征态必定是四重简并的。这与上一节得到的结果一致。

在 $U\gg g$ 的极限下，所有包含开弦的态都将具有量级为 $U$ 的能量。低能态只包含闭合弦网，并且满足
\begin{equation*}
F_{\pmb{i}}\big|_{\pmb{i}=\mathrm{even}}=1
\end{equation*}
不同的低能态由奇数方格上的 $F_{\pmb{i}}$ 标记，即 $F_{\pmb{i}}\big|_{\pmb{i}=\mathrm{odd}}=\pm1$。特别地，基态由下式给出
\begin{equation*}
F_{\pmb{i}}\big|_{\pmb{i}=\mathrm{odd}}=1.
\end{equation*}

这样的基态具有闭合弦凝聚。这是因为所有闭合弦算符 $W(C_{\mathrm{close}})$ 都与 $H_U+H_g$ 对易。于是，$H_U+H_g$ 的基态 $|\Psi_0\rangle$ 是 $W(C_{\mathrm{close}})$ 的本征态，并且满足
\begin{equation*}
\langle\Psi_0|W(C_{\mathrm{close}})|\Psi_0\rangle=1.
\end{equation*}

上式还意味着基态 $|\Psi_0\rangle$ 是各种闭合弦网组态的等权叠加。

基态之上的低能激发可以通过把某些奇数方格上的 $F_{\pmb{i}}$ 从 1 翻转到 $-1$ 得到。若把奇数方格上的 $F_{\pmb{i}}$ 视为 Z$_2$ 规范理论中的通量，我们发现模型的低能部分与 Z$_2$ 格点规范理论完全相同——至少对无限系统如此。这表明我们模型的低能有效理论是 Z$_2$ 格点规范理论。

然而，有人可能会反对这一结果，指出我们模型的低能部分也同样与在每个奇数方格上放一个自旋的伊辛模型相同，因此低能有效理论应当是伊辛模型。我们想指出，尽管对无限系统而言我们模型的低能部分与伊辛模型相同，但对有限系统而言，我们模型的低能部分不同于伊辛模型。例如，在具有周期性边界条件的有限偶 $\times$ 偶格子上，我们模型的基态有四重简并（Kitaev, 2003; Wen, 2003c）。伊辛模型没有这种简并，而 Z$_2$ 规范理论有这种简并。此外，我们的模型包含一种可以认定为 Z$_2$ 电荷（Z$_2$ charge）的激发。因此，我们模型的低能有效理论是 Z$_2$ 格点规范理论，而不是伊辛模型。奇数方格上 $F_{\pmb{i}}=-1$ 的激发可以视为 Z$_2$ 格点规范理论中的 Z$_2$ 涡旋激发。

为理解 Z$_2$ 电荷激发，我们注意到，在弦凝聚态中，弦不耗费能量，并且是不可观测的。把闭合弦算符作用到基态上，得到的仍是基态自身。由于弦不可观测，一段开弦的行为就像两个独立的粒子。开弦的每个端点对应偶数方格上的一个粒子。产生这样一个粒子所需的能量为 $U$ 的量级。现在，让我们考虑这样一个粒子围绕四个最近邻偶数方格的跳跃（见图 10.2）。每一步跳跃由一小段弦算符产生。我们看到，四个跳跃振幅的乘积由这四个偶数方格正中央那个奇数方格上 $\hat F_{\pmb{i}}$ 的本征值给出。这正是电荷与通量之间的关系。因此，若把奇数方格上的 $F_{\pmb{i}}$ 认定为 Z$_2$ 通量，则偶数方格上弦的端点就对应于 Z$_2$ 电荷。由于闭合弦凝聚，开弦的端点不被禁闭，彼此之间只有短程相互作用。因此，Z$_2$ 电荷的行为就像不挂着弦的点粒子。

%FIG{10.2}: {Z$_2$ 电荷围绕四个最近邻偶数方格的一次跳跃。}

\subsection{三种类型的弦与演生费米子}

\begin{keypoints}\item 弦的端点是规范荷。某些弦的端点还可以是费米子。\end{keypoints}

下面我们将研究几种不同类型的弦。为避免混淆，我们把上面讨论的弦称为 T1 弦。T1 弦连接偶数方格（见图 10.3(a)）。

正如 Z$_2$ 电荷一样，一对 Z$_2$ 涡旋也由开弦算符产生。由于 Z$_2$ 涡旋对应于奇数方格上被翻转的 $\hat F_{\pmb{i}}$，产生 Z$_2$ 涡旋的开弦算符同样由式 (10.3.1) 给出，只是现在乘积遍及一条连接\emph{奇数}方格的弦。我们把这样的弦称为 T2 弦。

我们想指出，T2 弦的参考态（即无弦态）不同于 T1 弦的参考态。无 T2 弦态由 $|\tilde{0}\rangle$ 给出，其自旋在偶数格点指向 $y$ 方向、在奇数格点指向 $x$ 方向。由于 T1 弦与 T2 弦有不同的参考态，我们不能同时拥有 T1 弦和 T2 弦的稀薄气体。容易验证，T2 弦算符也与 $H_J+H_g$ 对易。因此，基态 $|\Psi_0\rangle$ 除了 T1 闭合弦的凝聚之外，还有 T2 闭合弦的凝聚。

%FIG{10.3}: {三种类型的弦：(a) T1 弦；(b) T2 弦；(c) T3 弦。}

Z$_2$ 涡旋的跳跃由一段短的 T2 开弦诱导。由于 T2 开弦算符彼此都对易，Z$_2$ 涡旋的行为像玻色子。类似地，Z$_2$ 电荷的行为也像玻色子。然而，T1 开弦算符与 T2 开弦算符不对易。结果是，T1 弦的端点与 T2 弦的端点具有非平凡的互统计。鉴于我们已经证明，让 Z$_2$ 电荷绕 Z$_2$ 涡旋运动一周会产生相位 $\pi$，Z$_2$ 电荷与 Z$_2$ 涡旋具有半子型（semionic）互统计。

T3 弦定义为 T1 弦与 T2 弦的束缚态（见图 10.3(c)）。T3 弦算符由一个 T1 弦算符与一个 T2 弦算符的乘积给出，具有如下形式
\begin{equation*}
W(C)=\prod_n\sigma_{\pmb{i}_n}^{l_n}
\end{equation*}
其中 $C$ 是一条连接相邻链（link）中点的弦，$\pmb{i}_n$ 是弦上的格点。这里，若弦在格点 $\pmb{i}_n$ 处不拐弯，则 $l_n=z$；若弦在格点 $\pmb{i}_n$ 处拐弯，则 $l_n=x$ 或 $y$；若拐弯构成右上角或左下角，则 $l_n=x$；若拐弯构成右下角或左上角，则 $l_n=y$（见图 10.3(c)）。基态也有 T3 闭合弦的凝聚。T3 弦的端点是由一条链两侧两个方格上的一个 Z$_2$ 电荷与一个 Z$_2$ 涡旋形成的束缚态（即链两侧的 $F_{\pmb{i}}=-1$）。这样的束缚态是费米子（见 7.1.2 节与图 7.8）。所以费米子生活在链上。有趣的是，我们模型中的弦网凝聚直接导致 Z$_2$ 规范结构和三种新类型的准粒子，即 Z$_2$ 电荷、Z$_2$ 涡旋和费米子。费米子作为开 T3 弦的端点，从我们这个纯玻色模型中演生出来！

由于 T1 弦的端点是 Z$_2$ 电荷，T1 弦可以视为 Z$_2$“电”通量的弦。类似地，T2 弦可以视为 Z$_2$“磁”通量的弦。

\section{用投影对称群对不同的弦网凝聚分类}

\subsection{四类弦网凝聚}

\begin{keypoints}\item 不同的弦网凝聚给出不同的量子有序相。\end{keypoints}

我们已经看到，当 $U>0$、$g>0$、$J=0$ 时，我们模型 $H_g+H_U+H_J$ 的基态由下式给出
\begin{equation*}
F_{\pmb{i}}\big|_{\pmb{i}=\mathrm{even}}=1,\qquad F_{\pmb{i}}\big|_{\pmb{i}=\mathrm{odd}}=1.
\end{equation*}
我们将把这样的相称为 Z$_2$ 相，以强调其低能 Z$_2$ 规范结构。在 Z$_2$ 相中，T1 弦算符 $W_1(C_1)$ 与 T2 弦算符 $W_2(C_2)$ 有如下期望值：
\begin{equation*}
\langle W_1(C_1)\rangle=1,\qquad \langle W_2(C_2)\rangle=1
\end{equation*}

当 $U>0$、$g<0$、$J=0$ 时，基态由下式给出
\begin{equation*}
F_{\pmb{i}}\big|_{\pmb{i}=\mathrm{even}}=1,\qquad F_{\pmb{i}}\big|_{\pmb{i}=\mathrm{odd}}=-1.
\end{equation*}
我们看到每个奇数方格都有 $\pi$ 通量。我们将把这样的相称为 Z$_2$ 通量相。T1 弦算符与 T2 弦算符有如下期望值：
\begin{equation*}
\langle W_1(C_1)\rangle=(-)^{N_{\mathrm{odd}}},\qquad \langle W_2(C_2)\rangle=1
\end{equation*}
其中 $N_{\mathrm{odd}}$ 是被 T1 弦 $C_1$ 包围的奇数方格的数目。

当 $U<0$、$g>0$、$J=0$ 时，基态由下式给出
\begin{equation*}
F_{\pmb{i}}\big|_{\pmb{i}=\mathrm{even}}=-1,\qquad F_{\pmb{i}}\big|_{\pmb{i}=\mathrm{odd}}=1.
\end{equation*}
每个偶数方格上有一个 Z$_2$ 电荷。我们将把这样的相称为 Z$_2$ 电荷相。T1 弦算符与 T2 弦算符有如下期望值：
\begin{equation*}
\langle W_1(C_1)\rangle=1,\qquad \langle W_2(C_2)\rangle=(-)^{N_{\mathrm{even}}}
\end{equation*}
其中 $N_{\mathrm{even}}$ 是被 T2 弦 $C_2$ 包围的偶数方格的数目。注意，Z$_2$ 通量相与 Z$_2$ 电荷相仅相差一个格子平移，本质上是同一个相。

当 $U<0$、$g<0$、$J=0$ 时，基态变为
\begin{equation*}
F_{\pmb{i}}\big|_{\pmb{i}=\mathrm{even}}=-1,\qquad F_{\pmb{i}}\big|_{\pmb{i}=\mathrm{odd}}=-1.
\end{equation*}
每个偶数方格上有一个 Z$_2$ 电荷，且每个奇数方格有 $\pi$ 通量。我们将把这样的相称为 Z$_2$ 通量--电荷相。T1 弦算符与 T2 弦算符有如下期望值：
\begin{equation*}
\langle W_1(C_1)\rangle=(-)^{N_{\mathrm{odd}}},\qquad \langle W_2(C_2)\rangle=(-)^{N_{\mathrm{even}}}
\end{equation*}

当 $U=g=0$、$J\neq0$ 时，模型同样精确可解。基态是一个简单的自旋极化态（见图 10.1）。

从 $H_U+H_g+H_J$ 模型的这些精确可解极限出发，我们提出一幅相图，草图示于图 10.4。该相图包含四个不同的弦网凝聚相和一个无弦凝聚的相。所有的相都有相同的对称性，仅凭它们不同的量子序相互区分。

\subsection{投影对称群与凝聚弦的端点}

\begin{keypoints}\item PSG 所描述的投影对称性，正是凝聚弦端点有效理论的对称性。\end{keypoints}

从各不相同的 $\langle W_1(C_1)\rangle$ 与 $\langle W_2(C_2)\rangle$ 我们看到，前四个相具有不同的弦网凝聚。然而，它们都有相同的对称性。这就引出一个问题：在没有对称性破缺的情况下，我们怎么知道上述四个相的确是不同的相？我们怎么知道不可能在不经历相变的情况下把一个弦网凝聚态变成另一个弦网凝聚态？

%FIG{10.4}: {$H=H_U+H_g+H_J$ 模型的建议相图。这里假定 $J$ 为正。四个弦网凝聚相由一对 PSG（$PSG_{\mathrm{charge}}$，$PSG_{\mathrm{vortex}}$）表征。SP 标记自旋极化相。}

下面我们将证明，不同的弦网凝聚可以用不同的 PSG 来描述（正如不同的对称性破缺序可以用基态的不同对称群来描述）。在第 9 章中，不同的量子序是通过它们不同的 PSG 引入的（Wen, 2002b,c）。那里证明了 PSG 是量子相的普适性质，只有相变才能改变 PSG。因此，不同的弦网凝聚态由不同的 PSG 表征，这一事实表明这些不同的弦网凝聚态属于不同的量子相。

弦网凝聚与 PSG 之间的联系也使我们能够把弦网凝聚与第 9 章引入的量子序联系起来。事实上，第 9 章讨论的投影构造可以视为构造弦网凝聚态的一种方法。

为看清 PSG 与弦网凝聚之间的联系，我们注意到，当闭合弦网凝聚时，开弦的端点表现得像独立的粒子。让我们考虑描述 T1 弦两个端点的双粒子态 $|\pmb{p}_1\pmb{p}_2\rangle$。注意，T1 弦的端点（因而 Z$_2$ 电荷）只生活在偶数方格上。因此，$\pmb{p}_1$ 与 $\pmb{p}_2$ 只标记偶数方格。对于我们的模型 $H_U+H_g$，$|\pmb{p}_1\pmb{p}_2\rangle$ 是能量本征态，而且 Z$_2$ 电荷不跳跃。这里我们想往哈密顿量中加入一项
\begin{equation*}
H_t=t\sum_{\pmb{i}}(\sigma_{\pmb{i}}^{x}+\sigma_{\pmb{i}}^{y})+t'\sum_{\pmb{i}}\sigma_{\pmb{i}}^{z}
\end{equation*}
$t$ 项 $t\sum_{\pmb{i}}(\sigma_{\pmb{i}}^{x}+\sigma_{\pmb{i}}^{y})$ 使 Z$_2$ 电荷在偶数方格之间跳跃，跳跃幅为 $t$ 的量级。两个 Z$_2$ 电荷的动力学由双粒子希尔伯特空间中的如下有效哈密顿量描述：
\begin{equation*}
H=H(\pmb{p}_1)+H(\pmb{p}_2)
\end{equation*}
其中 $H(\pmb{p}_1)$ 描述第一个粒子 $\pmb{p}_1$ 的跳跃，$H(\pmb{p}_2)$ 描述第二个粒子 $\pmb{p}_2$ 的跳跃。现在我们可以定义弦网凝聚态的 PSG 了：PSG 就是跳跃哈密顿量 $H(\pmb{p})$ 的对称群。

由于底层模型 $H_U+H_g+H_t$ 具有平移对称性，我们可能天真地预期 Z$_2$ 电荷的跳跃哈密顿量 $H(\pmb{p})$ 在 $\pmb{x}+\pmb{y}$ 与 $\pmb{x}-\pmb{y}$ 方向也具有平移对称性：
\begin{align*}
H(\pmb{p})&=T_{xy}^{\dagger}H(\pmb{p})T_{xy}, & T_{xy}|\pmb{p}\rangle&=|\pmb{p}+\pmb{x}+\pmb{y}\rangle\\
H(\pmb{p})&=T_{x\bar{y}}^{\dagger}H(\pmb{p})T_{x\bar{y}}, & T_{x\bar{y}}|\pmb{p}\rangle&=|\pmb{p}+\pmb{x}-\pmb{y}\rangle\tag{10.4.1}
\end{align*}
如果上式成立，则意味着 PSG 与平移对称群相同。

结果表明，式 (10.4.1) 过强。即使 $H(\pmb{p})$ 不满足式 (10.4.1)，底层自旋模型仍然可以具有平移对称性。然而，$H(\pmb{p})$ 可能的对称群（即各个 PSG）受到底层自旋模型平移对称性的很强约束。下面我们将解释，为什么 PSG 可以不同于物理自旋模型的对称群，以及为与底层自旋模型的平移对称性相一致，PSG 必须满足什么条件。

我们注意到，弦总有两个端点。因此，一个物理态总是拥有偶数个 Z$_2$ 电荷。平移作用在双粒子态上为
\begin{align*}
T_{xy}^{(2)}|\pmb{p}_1,\pmb{p}_2\rangle&=|\pmb{p}_1+\pmb{x}+\pmb{y},\ \pmb{p}_2+\pmb{x}+\pmb{y}\rangle\\
T_{x\bar{y}}^{(2)}|\pmb{p}_1,\pmb{p}_2\rangle&=|\pmb{p}_1+\pmb{x}-\pmb{y},\ \pmb{p}_2+\pmb{x}-\pmb{y}\rangle
\end{align*}
这里 $T_{xy}^{(2)}$ 与 $T_{x\bar{y}}^{(2)}$ 满足平移代数
\begin{equation*}
T_{xy}^{(2)}T_{x\bar{y}}^{(2)}=T_{x\bar{y}}^{(2)}T_{xy}^{(2)}\tag{10.4.2}
\end{equation*}

我们注意到，$T_{xy}^{(2)}$ 与 $T_{x\bar{y}}^{(2)}$ 是单粒子态上平移算符的直积。因此，在某种意义上，单粒子平移是双粒子平移的平方根。双粒子平移代数对单粒子平移代数施加了很强的约束。

单粒子平移最一般的形式由 $T_{xy}G_{xy}$ 与 $T_{x\bar{y}}G_{x\bar{y}}$ 给出，其中算符 $T_{xy,x\bar{y}}$ 与 $G_{xy,x\bar{y}}$ 的作用定义为
\begin{align*}
T_{xy}|\pmb{p}\rangle&=|\pmb{p}+\pmb{x}+\pmb{y}\rangle, & T_{x\bar{y}}|\pmb{p}\rangle&=|\pmb{p}+\pmb{x}-\pmb{y}\rangle,\\
G_{xy}|\pmb{p}\rangle&=\mathrm{e}^{\phi_{xy}(\pmb{p})}|\pmb{p}\rangle, & G_{x\bar{y}}|\pmb{p}\rangle&=\mathrm{e}^{\phi_{x\bar{y}}(\pmb{p})}|\pmb{p}\rangle.
\end{align*}
为了使直积 $T_{xy}^{(2)}=T_{xy}G_{xy}\otimes T_{xy}G_{xy}$ 与 $T_{x\bar{y}}^{(2)}=T_{x\bar{y}}G_{x\bar{y}}\otimes T_{x\bar{y}}G_{x\bar{y}}$ 复现平移代数 (10.4.2)，我们只要求 $T_{xy}G_{xy}$ 与 $T_{x\bar{y}}G_{x\bar{y}}$ 满足
\begin{equation*}
T_{xy}G_{xy}T_{x\bar{y}}G_{x\bar{y}}=T_{x\bar{y}}G_{x\bar{y}}T_{xy}G_{xy}\tag{10.4.3}
\end{equation*}
或
\begin{equation*}
T_{xy}G_{xy}T_{x\bar{y}}G_{x\bar{y}}=-T_{x\bar{y}}G_{x\bar{y}}T_{xy}G_{xy}\tag{10.4.4}
\end{equation*}

算符 $T_{xy}G_{xy}$ 与 $T_{x\bar{y}}G_{x\bar{y}}$ 生成一个群，这个群正是 PSG。两个不同的代数 (10.4.3) 与 (10.4.4) 生成两个不同的 PSG，二者都与作用在双粒子态上的平移群相一致。我们把由式 (10.4.3) 生成的 PSG 称为 Z2a PSG，把由式 (10.4.4) 生成的 PSG 称为 Z2b PSG。

让我们回顾一下 PSG 的一般定义。PSG 是一个群，它是对称群（SG）的扩张，即 PSG 包含一个正规子群（称为不变规范群或 IGG），使得
\begin{equation*}
PSG/IGG=SG
\end{equation*}

在我们的情形中，SG 是平移群 $SG=\{1,T_{xy}^{(2)},T_{x\bar{y}}^{(2)},\dots\}$。IGG 由单粒子态上满足 $G_0\otimes G_0=1$ 的变换 $G_0$ 构成。我们发现 IGG 由下式生成：
\begin{equation*}
G_0|\pmb{p}\rangle=-|\pmb{p}\rangle
\end{equation*}

单粒子态上的这三个变换 $G_0$、$T_{xy}G_{xy}$ 与 $T_{x\bar{y}}G_{x\bar{y}}$ 生成 Z2a 或 Z2b PSG。

我们现在看到，底层平移对称性并不要求单粒子跳跃哈密顿量 $H(\pmb{p})$ 具有平移对称性，只要求 $H(\pmb{p})$ 在 Z2a PSG 或 Z2b PSG 下不变。当 $H(\pmb{p})$ 在 Z2a PSG 下不变时，跳跃哈密顿量具有通常的平移对称性。当 $H(\pmb{p})$ 在 Z2b PSG 下不变时，跳跃哈密顿量具有磁平移对称性，描述磁场中每个奇数方格带 $\pi$ 通量的跳跃。

\subsection{投影对称群对不同的弦网凝聚的分类}

\begin{keypoints}\item 不同的弦网凝聚有不同的 PSG。由于 PSG 是量子相的普适性质，我们发现不同的弦网凝聚对应于不同的相。\end{keypoints}

在了解了弦端点的跳跃哈密顿量可能具有的 PSG 之后，我们现在可以计算实际的 PSG 了。让我们考虑模型 $H_U+H_g+H_t$ 的两个基态：一个具有 $F_{\pmb{i}}\big|_{\pmb{i}=\mathrm{odd}}=1$（对应 $g>0$），另一个具有 $F_{\pmb{i}}\big|_{\pmb{i}=\mathrm{odd}}=-1$（对应 $g<0$）。两个基态在 $\pmb{x}+\pmb{y}$ 与 $\pmb{x}-\pmb{y}$ 方向都有相同的平移对称性。然而，相应的单粒子跳跃哈密顿量 $H(\pmb{p})$ 却有不同的对称性。对于 $F_{\pmb{i}}\big|_{\pmb{i}=\mathrm{odd}}=1$ 的态，奇数方格没有通量，$H(\pmb{p})$ 具有通常的平移对称性，在 Z2a PSG 下不变。对于 $F_{\pmb{i}}\big|_{\pmb{i}=\mathrm{odd}}=-1$ 的态（译注：原文此处印作“$F_{\pmb{i}}|_{\pmb{i}=\mathrm{odd}}=1$”，按上下文应为 $-1$），奇数方格有 $\pi$ 通量，$H(\pmb{p})$ 具有磁平移对称性，其 PSG 为 Z2b PSG。因此，$F_{\pmb{i}}\big|_{\pmb{i}=\mathrm{odd}}=1$ 的态与 $F_{\pmb{i}}\big|_{\pmb{i}=\mathrm{odd}}=-1$ 的态尽管对称性相同，却有不同的序。两个态中不同的量子序可以用它们不同的 PSG 来表征。

如 10.4.1 小节所述，上述两个态有不同的弦网凝聚。对 $\hat F_{\pmb{i}}\big|_{\pmb{i}=\mathrm{odd}}=1$ 的态，T1 闭合弦算符的平均值为 $\langle W(C)\rangle=1$；而对 $\hat F_{\pmb{i}}\big|_{\pmb{i}=\mathrm{odd}}=-1$ 的态，T1 闭合弦算符的平均值为 $\langle W(C)\rangle=(-)^{N_{\mathrm{odd}}(C)}$，其中 $N_{\mathrm{odd}}(C)$ 是被 T1 闭合弦 $C$ 包围的奇数方格的数目。因此，我们也可以说，不同的弦网凝聚可以用不同的 PSG 来表征。

在上面的讨论中，我们只在一些特定的精确可解模型中表明了 PSG 可以表征不同的弦网凝聚态。事实上，在一般模型中，不同的 PSG 确实表征不同的量子相。这是因为，正如 9.9.1 小节所证明的，PSG 是相的普适性质（Wen, 2002c）。弦网凝聚态的 PSG 对物理模型的任何不改变模型对称性的局域微扰都是稳健的。

上述讨论同样适用于 Z$_2$ 涡旋与 T2 弦。因此，我们模型中的量子序由一对 PSG（$PSG_{\mathrm{charge}}$，$PSG_{\mathrm{vortex}}$）描述，一个对应 Z$_2$ 电荷，一个对应 Z$_2$ 涡旋。PSG 对（$PSG_{\mathrm{charge}}$，$PSG_{\mathrm{vortex}}$）使我们能够区分模型 $H=H_U+H_g+H_t$ 的四个不同的弦网凝聚态（见图 10.5）。

\subsection{T3 弦端点的投影对称群}

\begin{keypoints}\item 对于具有几种不同类型弦凝聚的态，不同凝聚弦的端点可以有不同的投影对称性和不同的 PSG。\end{keypoints}

%FIG{10.5}: {$H=H_U+H_g+H_t$ 模型的建议相图。这里假定 $t=t'$ 为正。四个弦网凝聚相由 PSG 对（$PSG_{\mathrm{charge}}$，$PSG_{\mathrm{vortex}}$）表征。SP 标记均匀自旋极化相。所有的相都有相同的对称性，并且没有任何相变改变对称性。}

\begin{keypoints}\item 9.4.2 小节引入的 PSG 对应于 T3 弦的 PSG。\end{keypoints}

在本小节中，我们将假设模型中 $U=g=V$：
\begin{equation*}
H_U+H_g+H_t=H_t-V\sum_{\pmb{i}}\hat F_{\pmb{i}}\tag{10.4.5}
\end{equation*}

新模型具有由 $\pmb{i}\to\pmb{i}+\pmb{x}$ 与 $\pmb{i}\to\pmb{i}+\pmb{y}$ 生成的更大的平移对称性（见图 10.5）。注意到，当 $H_t=0$ 时（译注：原文此处印作“$H_l=0$”，应为 $H_t$），新模型与 10.2.1 小节研究过的自旋哈密顿量 (10.2.6) 相同。凝聚的 T1 弦与 T2 弦的 PSG 已在上一小节研究过，这里我们想讨论 T3 弦的 PSG。由于 T3 弦的端点生活在链上，相应的单粒子跳跃哈密顿量 $H_f(\pmb{l})$ 描述费米子在链之间的跳跃。显然，$H_f(\pmb{l})$ 的对称群（即 PSG）可以不同于 $H(\pmb{p})$ 的对称群。

让我们把费米子跳跃 $\pmb{l}\to\pmb{l}+\pmb{x}$ 定义为两步跳跃 $\pmb{l}\to\pmb{l}+\frac{\pmb{x}}{2}-\frac{\pmb{y}}{2}\to\pmb{l}+\pmb{x}$ 的组合（见图 10.6）。跳跃 $\pmb{l}\to\pmb{l}+\frac{\pmb{x}}{2}-\frac{\pmb{y}}{2}$ 由 $\sigma_{\pmb{l}+\frac{\pmb{x}}{2}}^{x}$ 产生，跳跃 $\pmb{l}+\frac{\pmb{x}}{2}-\frac{\pmb{y}}{2}\to\pmb{l}+\pmb{x}$ 由 $\sigma_{\pmb{l}+\frac{\pmb{x}}{2}}^{y}$ 产生。因此，跳跃 $\pmb{l}\to\pmb{l}+\pmb{x}$ 由 $\sigma_{\pmb{l}+\frac{\pmb{x}}{2}}^{y}\sigma_{\pmb{l}+\frac{\pmb{x}}{2}}^{x}$ 产生。类似地，我们把费米子跳跃 $\pmb{l}\to\pmb{l}+\pmb{y}$ 定义为两步跳跃 $\pmb{l}\to\pmb{l}+\frac{\pmb{x}}{2}+\frac{\pmb{y}}{2}\to\pmb{l}+\pmb{y}$，它由 $\sigma_{\pmb{l}+\frac{\pmb{x}}{2}+\pmb{y}}^{x}\sigma_{\pmb{l}+\frac{\pmb{x}}{2}}^{y}$ 产生。
%FIG{10.6}: {费米子绕一个方格和绕一个格点的跳跃。}
在这样的定义下，费米子绕方格 $\pmb{l}\to\pmb{l}+\pmb{x}\to\pmb{l}+\pmb{x}+\pmb{y}\to\pmb{l}+\pmb{y}\to\pmb{l}$ 跳跃一周由下式产生（见图 10.6）：$(\sigma_1^{y}\sigma_4^{x})(\sigma_4^{x}\sigma_4^{y})(\sigma_3^{x}\sigma_2^{y})(\sigma_1^{y}\sigma_1^{x})=\sigma_4^{y}\sigma_3^{x}\sigma_2^{y}\sigma_1^{x}=\hat F_1$，或者 $(\sigma_5^{z})(\sigma_5^{y}\sigma_8^{x})(\sigma^{z})(\sigma_7^{x}\sigma_6^{y})=-\sigma_5^{x}\sigma_8^{y}\sigma_7^{x}\sigma_6^{y}=-\hat F_5$。基态具有 $F_{\pmb{i}}=\mathrm{sgn}(V)$。然而，由于 T3 弦终止在链 $5\to6$ 上，我们有 $-\hat F_5=\hat F_{\pmb{i}}=\mathrm{sgn}(V)$；另一方面，我们有 $\hat F_1=\hat F_{\pmb{i}}=\mathrm{sgn}(V)$。因此，费米子绕 1234 与 5678 两个方格跳跃的总振幅由式 (10.4.5) 中 $V$ 的符号给出。当 $\mathrm{sgn}(V)=1$ 时，费米子看不到通量；而当 $\mathrm{sgn}(V)=-1$ 时，费米子看到每个方格有 $\pi$ 通量。结果是，$H_f(\pmb{l})$ 平移对称性的生成元 $(T_xG_x,\,T_yG_y)$ 满足磁平移代数（见式 (7.2.3)）
\begin{equation*}
(T_yG_y)^{-1}(T_xG_x)^{-1}T_yG_yT_xG_x=\mathrm{sgn}(V).\tag{10.4.6}
\end{equation*}

除了在 $(T_xG_x,\,T_yG_y)$ 下不变之外，$H_f(\pmb{l})$ 在 $G_0$ 下也不变：
\begin{equation*}
G_0|\pmb{l}\rangle=-|\pmb{l}\rangle
\end{equation*}
于是，$(G_0,\,T_xG_x,\,T_yG_y)$ 生成 $H_f(\pmb{l})$ 的对称群——费米子 PSG。

当 $\mathrm{sgn}(V)=1$ 时，$T_xG_x$ 与 $T_yG_y$ 满足平移代数，相应的费米子 PSG 是 Z2A PSG。当 $\mathrm{sgn}(V)=-1$ 时，$T_xG_x$ 与 $T_yG_y$ 满足每个方格带 $\pi$ 通量的磁平移代数，相应的费米子 PSG 是 Z2B PSG。我们看到，基态的量子序也可以用费米子 PSG 来表征。$V<0$ 基态的量子序由 Z2A PSG 表征，而 $V>0$ 基态的量子序由 Z2B PSG 表征。

在 10.2 节中，自旋 1/2 模型 (10.4.5)（取 $t=t'=0$）曾通过把自旋拆成两个费米子、用从属玻色子方法求解。那里证明了，$V>0$ 与 $V<0$ 两种态的费米子跳跃哈密顿量具有不同的对称性，或者说在不同的 PSG 下不变。不同的 PSG 意味着 $V>0$ 与 $V<0$ 基态中不同的量子序。10.2 节中针对 $V>0$ 与 $V<0$ 相得到的 PSG 与
